\section*{الفصل \ref{ch:b3:spectral} --- المؤثرات المتراصة
والمبرهنة الطيفية}

\begin{solution}{exo:b3:spectral:1}
(a) المجموعة $T(B)$ جزء محدود من الفضاء المنتهي البُعد $\operatorname{im}T$:
فهي متراصة نسبيًّا بهاينه--بوريل (\cref{cor:b3:topology:heineborel}، منقولةً \omterm{def:b3:topology:continuity}{بتماثل
طوبولوجي} خطي مع $\R^n$).
(b) وكون $I$ \omterm{def:b3:spectral:compact}{متراصًّا} يعني تراص الكرة الواحدية المغلقة، وهو
يقع، حسب مبرهنة ريس (السنة الجامعية 2)، في البُعد المنتهي
بالضبط. ولو كان \omterm{def:b3:spectral:compact}{لمؤثر متراص} $T$ مقلوبٌ محدود $T^{-1}$، لكان
$I = T^{-1}T$ \omterm{def:b3:spectral:compact}{متراصًّا} (\cref{prop:b3:spectral:ideal}): وهو مستحيل في البُعد
اللانهائي.
\end{solution}

\begin{solution}{exo:b3:spectral:2}
(a) $\norm{Tx}^2 = \sum\abs{\lambda_n}^2\abs{x_n}^2 \leq
\sup\abs{\lambda_n}^2\norm x^2$، مع مقاربة التساوي على المتجهات $e_n$ التي تحقق
السوپريموم.
(b) ($\Leftarrow$) للبتور $T_N$ (بالاحتفاظ بالأدلّة $n \leq N$
والإلغاء بعدها) رتبةٌ منتهية و $\vertiii{T - T_N} =
\sup_{n>N}\abs{\lambda_n} \to 0$: فهي متراصة حسب
\cref{prop:b3:spectral:ideal}. ($\Rightarrow$) وإذا كان $\abs{\lambda_{n_k}} \geq \delta > 0$ على امتداد متتالية
جزئية: $\norm{Te_{n_k} - Te_{n_l}}^2 = \abs{\lambda_{n_k}}^2 +
\abs{\lambda_{n_l}}^2 \geq 2\delta^2$: فلا توجد أي متتالية جزئية متقاربة من
$(Te_{n_k})$.
(c) المؤثر $T^*$ هو القطري بالمعاملات $(\bar\lambda_n)$: فهو \omterm{def:b3:spectral:selfadjoint}{ذاتي المرافقة} إذا
وفقط إذا كانت جميع $\lambda_n \in \R$. عندئذٍ يكون الأساس المعياري $(e_n)$
أساسًا متعامدًا متجانسًا من المتجهات الذاتية، بقيم ذاتية
$\lambda_n \to
0$: أي المبرهنة الطيفية حرفيًّا.
\end{solution}

\begin{solution}{exo:b3:spectral:3}
لتكن $(x_n)$ محدودة. عندئذٍ تكون $(Bx_n)$ محدودة ($\vertiii B < \infty$)؛
ويستخرج تراص $S$ المتتاليةَ $SBx_{n_k} \to y$؛ ويعطي اتصال $A$ أن $ASBx_{n_k} \to
Ay$:
أي إن $ASB$ متراص. وإذا كان $ST = TS = I$ مع $T$ متراص وكان
$\dim H = \infty$: لكان $I = ST$ \omterm{def:b3:spectral:compact}{متراصًّا}، وهو يناقض \cref{exo:b3:spectral:1}(b) ---
وهي العبارة نفسها منظورًا إليها من الجهة الأخرى.
\end{solution}

\begin{solution}{exo:b3:spectral:4}
(a) بكوشي--شوارتز في $y$: $\abs{T_kf(x)}^2 \leq
\bigl(\int\abs{k(x,y)}^2\dd y\bigr)\norm f_2^2$؛ ثم نكامل في $x$ (بتونيلي):
$\norm{T_kf}_2 \leq \norm k_{L^2(\square)}
\norm f_2$.

(b) العائلة $e_{mn}(x,y) = e_m(x)\overline{e_n(y)}$ متعامدة متجانسة في $L^2(\intcc01^2)$ (إذ تفصل تونيلي
التكاملَ المزدوج). وأما الشمول: فإذا كانت $h \perp$ جميع
$e_{mn}$، فإنه من أجل كل $m$ تكون الدالة $y \mapsto \int h(x,y)\overline{e_m(x)}\dd
x$ (وهي في $L^2$
بكوشي--شوارتز وتونيلي) متعامدةً مع كل $\overline{e_n}$ ---
والمرافقات $(\overline{e_n})$ تشكّل \omterm{def:b3:hilbert:onb}{أساسًا هيلبرتيًّا} كلما فعلت $(e_n)$
ذلك (إذ المرافقة تقابل متقايس للفضاء $L^2$ يحفظ التعامد
والشمول) --- ومنه فهي $0$ في كل مكان تقريبًا؛ ثم من أجل $y$ في
كل مكان تقريبًا، $h(\cdot, y) \perp$ كل $e_m$: أي $h(\cdot, y) = 0$ في كل مكان تقريبًا:
ومنه $h = 0$ (بتونيلي). فتكون $(e_{mn})$ \omterm{def:b3:hilbert:onb}{أساسًا هيلبرتيًّا}؛
وننشر $k = \sum c_{mn}e_{mn}$. ويعطي البتر $k_N$ (بالأدلّة $\leq N$) مؤثرًا
$T_{k_N}$ منتهي الرتبة (مجاله في $\operatorname{Vect}(e_1, \dots,
e_N)$)، وحسب (a)،
\[
\vertiii{T_k - T_{k_N}} \leq \norm{k - k_N}_{L^2} \to 0 :
\]
فيكون $T_k$ نهايةً بالمعيار لمؤثرات منتهية الرتبة: أي \omterm{def:b3:spectral:compact}{متراصًّا}
(\cref{prop:b3:spectral:ideal}).
\end{solution}

\begin{solution}{exo:b3:spectral:5}
(a) $\lambda\norm x^2 = \langle x, Tx\rangle \geq 0$ على متجهة ذاتية. والصيغة هي \cref{prop:b3:spectral:sanorm} مع جميع
القيم $\langle x,
Tx\rangle \geq 0$: فالقيمة المطلقة زائدة.
(b) المقدار $(x, y) \mapsto \langle x, Ty\rangle$ شكلٌ نصف خطي إرميتي موجب (وقد يكون منحلًّا)؛
وبرهان كوشي--شوارتز المعتاد (بنشر $\langle x + t\eu^{\iu\theta}y,
T(x + t\eu^{\iu\theta}y)\rangle \geq 0$ وأخذ المميِّز) لا
يستعمل التعيين البتة.
\end{solution}

\begin{solution}{exo:b3:spectral:6}
(a) $\norm{Mf}_2 \leq \norm f_2$، وعلى $f_n = \sqrt n\,
\mathbf 1_{\intcc{1 - 1/n}1}$ (وهي متجهات واحدية)، $\norm{Mf_n}_2
\geq 1 - \frac1n$: ومنه
$\vertiii M = 1$؛ وهو \omterm{def:b3:spectral:selfadjoint}{ذاتي المرافقة} لأن عامل الضرب حقيقي. وأما
القيم الذاتية: فإن $xf(x) = \lambda f(x)$ في كل مكان تقريبًا يفرض $f = 0$ في كل
مكان تقريبًا خارج المجموعة المعدومة $\{x = \lambda\}$: أي $f =
0$ في $L^2$.
(b) وبالدوال $f_n$ نفسها: $\norm{Mf_n - f_n}_2 \leq \frac1n \to
0$، بينما $f_n \rightharpoonup 0$ (إذ من أجل
$g \in L^2$ مثبَّتة، $\abs{\langle g, f_n\rangle} \leq \norm{g\,\mathbf
1_{\intcc{1-1/n}1}}_2 \to 0$ بالتقارب المهيمن). فإذا كان
$Mf_{n_k} \to h$ بالمعيار، لكان $f_{n_k} \to h$، وهو يفرض $h = 0$
(بالنهاية الضعيفة) مع أن $\norm h = 1$: فلا توجد أي متتالية
جزئية متقاربة.
(c) في \cref{lem:b3:spectral:existence}، يستعمل التراصَ بالضبط الاستخراجُ
«$Tx_{n_k} \to y$»؛ ومن أجل $M$ تتركّز المتتاليات المعظِّمة
بجوار $x = 1$ وتتقارب صورها ضعيفًا إلى $0$ ولا تتقارب بالمعيار
أبدًا: فببساطة لا توجد المتجهة الذاتية عند قمة المجال العددي.
\end{solution}

\begin{solution}{exo:b3:spectral:7}
(a) $V = T_k$ مع $k(x, y) = \mathbf 1_{y < x} \in
L^2(\intcc01^2)$: فهو متراص (\cref{exo:b3:spectral:4}). ومرافقه مؤثر
النواة ذو النواة $\overline{k(y,
x)} = \mathbf 1_{y > x}$: أي $V^*f(x) = \int_x^1f$؛ و $V \neq V^*$ (بالاختبار
على $f = \mathbf 1$).
(b) إذا كان $Vf = \lambda f$ مع $\lambda \ne 0$: فإن $Vf$ متصلة
على $\intcc01$ (بالتقارب المهيمن في $\int_0^x f$)، ومنه تقبل
$f
= \frac1\lambda Vf$ ممثّلًا \omterm{def:b3:topology:continuity}{متصلًا}؛ عندئذٍ تكون $Vf$ من الصنف
$\mathcal C^1$ (بالمبرهنة الأساسية في التفاضل والتكامل من أجل
مقادير مكامَلة متصلة)، ومنه تكون $f$ من الصنف $\mathcal C^1$،
و $\lambda f' = f$ مع $f(0) = \frac1\lambda Vf(0) = 0$: أي $f =
C\eu^{x/\lambda}$ مع $C = f(0) = 0$.
(c) فالمؤثر $V$ متراص وليست له أي قيمة ذاتية البتة عدا $0$
ربما (إذ يفرض $Vf = 0$ أن $f = 0$ في كل مكان تقريبًا باشتقاق
التكامل --- ومنه ولا حتى $0$): أي إن الآلة الطيفية تتطلب فعلًا
ذاتية المرافقة، لا التراص وحده.
\end{solution}

\begin{solution}{exo:b3:spectral:8}
نكتب $x = \sum_ic_ie_i + z$ حيث $z \in \ker T$، ومنه $\langle x,
Tx\rangle = \sum_i\mu_i\abs{c_i}^2$. \emph{الحدّ
الأدنى}: على الكرة الواحدية للفضاء $V_k = \operatorname{Vect}(e_1, \dots,
e_k)$، $\langle x, Tx\rangle = \sum_{i\leq
k}\mu_i\abs{c_i}^2 \geq \mu_k$: ومنه يكون
الأعظم على $V$ للأصغر $\geq \mu_k$. \emph{الحدّ الأعلى}: ليكن
$\dim V = k$ ولتكن $W_k
= \overline{\operatorname{Vect}}(e_k, e_{k+1}, \dots) +
\ker T$، وبُعدها المتمم $k - 1$ (إذ متممها
المتعامد هو $V_{k-1}$)؛ فإن $V \cap W_k \neq \{0\}$ (لأن لتطبيق خطي $V \to
H/W_k \cong V_{k-1}$ رتبته
$\leq k - 1$ نواةً غير بديهية)، ولمتجهة واحدية $x \in V\cap W_k$ يكون
$\langle x,
Tx\rangle = \sum_{i \geq k}\mu_i\abs{c_i}^2 \leq \mu_k$: ومنه يكون الأصغر على $V$ $\leq \mu_k$.
وبضمّ ذلك: نجد الصيغة الأولى؛ والثانية تُبرهَن بالتناظر (باختبار
$W = W_k$؛ ومن أجل $W$ كيفية ذات بُعد متمم $k-1$، يعطي $W \cap V_k \neq 0$
متجهةً واحدية تحقق $\langle x, Tx\rangle \geq \mu_k$). وأما الرتابة: فينتقل $\langle x, Tx\rangle \leq \langle x,
Sx\rangle$ نقطةً
نقطة عبر $\max\min$.
\end{solution}

\begin{solution}{exo:b3:spectral:9}
نعيد كتابة $f - \lambda Tf = g$. ومن أجل $\lambda = 0$: يكون $f = g$، وهو
قابل للحل بكيفية وحيدة دائمًا. ومن أجل $\lambda \neq 0$: تصير
$(T - \frac1\lambda)f = -\frac g\lambda$، وحسب \omterm{thm:b3:spectral:fredholm}{بديل فريدهولم} (\cref{thm:b3:spectral:fredholm}) مع القيم الذاتية
$\lambda_n = \bigl((n + \frac12)\pi\bigr)^{-2}$ للمؤثر $T$ (\cref{ex:b3:spectral:minxy}): تكون قابلية الحل الوحيد من أجل
كل $g$ متحققةً إذا وفقط إذا كان $\frac1\lambda \neq \lambda_n$ من أجل كل $n$، أي
\[
\lambda \neq \Bigl(n + \tfrac12\Bigr)^2\pi^2
\qquad (n = 0, 1, 2, \dots).
\]
وعند قيمة استثنائية $\lambda = (n+\frac12)^2\pi^2$: توجد الحلول إذا وفقط إذا كان
$g \perp \sin\bigl((n{+}\frac12)\pi x\bigr)$، وتكون عندئذٍ وحيدة إلى غاية إضافة مضاعفات لذلك الجيب.
\end{solution}

\begin{solution}{exo:b3:spectral:10}
(a) $Sx = \lambda x$: بمقارنة الإحداثيات، $0 = \lambda
x_1$ و $x_n = \lambda x_{n+1}$؛
فإذا كان $\lambda \ne 0$ فإن $x_1 = 0$ وبالتراجع $x = 0$؛ وإذا
كان $\lambda = 0$، فإن $Sx = 0$ يفرض $x = 0$ (لأن $S$ متقايس).
فلا قيم ذاتية. وتُقرأ $S^*x =
\lambda x$ على الصورة $x_{n+1} = \lambda x_n$: أي $x =
x_1(1, \lambda, \lambda^2, \dots)$، وهي
في $\ell^2$ إذا وفقط إذا كان $\abs\lambda < 1$: أي قرص مفتوح كامل من القيم
الذاتية.
(b) $\norm{Se_n - Se_m} = \norm{e_{n+1} - e_{m+1}} = \sqrt2$: فليس لصورة المتتالية المحدودة $(e_n)$ أي متتالية
جزئية كوشية؛ وبالمثل $S^*e_{n+1} = e_n$. فلا هذا ولا ذاك متراص.
(c) من أجل المؤثرات المتراصة ذاتية المرافقة تلتقط القيم الذاتية
كل شيء (\cref{thm:b3:spectral:spectral})؛ وبإسقاط أي من الفرضيتين، قد تختفي القيم
الذاتية تمامًا ($S$، وفولتيرا) أو تملأ قرصًا ($S^*$): فالكائن
المتين هو الطيف $\{\lambda : T - \lambda I \text{ غير قابل للقلب}\}$، ونظريته تنتمي إلى مقرر لاحق.
\end{solution}

\begin{solution}{exo:b3:spectral:11}
(a) المؤثر $S$ هو المؤثر القطري بالمعاملات $\sqrt{\mu_n} \to 0$: فهو متراص
(\cref{exo:b3:spectral:2}(b)، منقولًا إلى الأساس $(e_n)$ مكمَّلًا بالنواة
$\ker T$ حيث $S = 0$)، \omterm{def:b3:spectral:selfadjoint}{وذاتي المرافقة} (لأنه قطري حقيقي)، وموجب
($\langle x, Sx\rangle =
\sum\sqrt{\mu_n}\abs{\langle e_n, x\rangle}^2$)، و $S^2 =
T$ حدًّا حدًّا.

(b) يتبادل $R$ مع $T = R^2$؛ ومن أجل متجهة ذاتية $x$ للمؤثر
$T$ بقيمة ذاتية $\mu$: $T(Rx) = RTx = \mu Rx$، ومنه يكون الفضاء الذاتي المنتهي
البُعد $E_\mu = \ker(T - \mu)$ مستقرًّا بالمؤثر $R$. وعلى $E_\mu$، يكون $R$
متناظرًا موجبًا مع $R^2
= \mu\,\mathrm{id}$: فتحقق قيمه الذاتية $\rho$ أن
$\rho^2
= \mu$ مع $\rho \geq 0$: ومنه تساوي كلها $\sqrt\mu$، ومؤثر
قابل للقُطرنة بقيمة ذاتية واحدة سلّميٌّ: أي $R = \sqrt\mu\,\mathrm{id}$ على
$E_\mu$. وعلى $\ker T$: $\norm
{Rx}^2 = \langle x, R^2x\rangle = \langle x, Tx\rangle = 0$. فيتوافق $R$ مع $S$ على
$\ker T$ وعلى كل فضاء ذاتي، والفضاء المولَّد المغلق بها هو $H$
(بالمبرهنة الطيفية): ومنه $R = S$.

(c) يفعل $\sqrt G$ بالمقدار $\frac1{n\pi}$ على $e_n =
\sqrt2\sin(n\pi x)$: فهو
مؤثر النواة ذو
\[
k(x, y) = \sum_{n\geq1}\frac{2\sin(n\pi x)\sin(n\pi
y)}{n\pi},
\]
والمتسلسلة متقاربة في $L^2(\intcc01^2)$ (بمعاملات $\frac1{n\pi} \in \ell^2$؛ وتعطي نوى
المجاميع الجزئية التقريبات المنتهية الرتبة). ولا حاجة إلى أي
صيغة مغلقة ابتدائية: فالجانب الطيفي \emph{هو} المؤثر.
\end{solution}

\begin{solution}{exo:b3:spectral:12}
(a) المؤثر $T^*T$ متراص (لأنه جداء مؤثر محدود \omterm{def:b3:spectral:compact}{ومؤثر متراص}،
\cref{exo:b3:spectral:3})، \omterm{def:b3:spectral:selfadjoint}{وذاتي المرافقة} ($(T^*T)^* = T^*T$)، وموجب ($\langle x, T^*Tx\rangle =
\norm{Tx}^2$). وأما
النواة: $T^*Tx = 0 \Rightarrow \norm{Tx}^2 =
\langle x, T^*Tx\rangle = 0 \Rightarrow Tx = 0$، وبالعكس: $\ker T^*T = \ker T$. وتوفّر المبرهنة الطيفية
العائلةَ المتعامدة المتجانسة $(e_n)$ مع $T^*Te_n = s_n^2e_n$ و $s_n > 0$،
المولِّدةَ للفضاء $(\ker T)^\perp$.

(b) $\langle f_m, f_n\rangle = \frac{\langle Te_m,
Te_n\rangle}{s_ms_n} = \frac{\langle e_m,
T^*Te_n\rangle}{s_ms_n} = \frac{s_n^2}{s_ms_n}\delta_{mn} =
\delta_{mn}$. وننشر $x = x_0 + \sum_n\langle e_n, x\rangle
e_n$ مع $x_0 \in \ker T$ (\omterm{thm:b3:hilbert:parseval}{ببارسيفال} في
الفضاء المولَّد المغلق مع النواة)؛ وبتطبيق $T$ المتصل:
\[
Tx = \sum_n\langle e_n, x\rangle\,Te_n =
\sum_ns_n\langle e_n, x\rangle\,f_n,
\]
والمتسلسلة متقاربة لأن مجاميعها الجزئية كوشية ($\norm{\sum_{N<n\leq M}s_n\langle e_n, x\rangle f_n}^2 =
\sum s_n^2\abs{\langle e_n, x\rangle}^2$،
مهيمَنًا عليها بالمقدار $\sup_{n>N}s_n^2\cdot\norm x^2$، و $s_n \to 0$).

(c) $\norm{Tx}^2 = \sum_ns_n^2\abs{\langle e_n, x\rangle}^2
\leq (\max s_n)^2\norm x^2$، وهو مبلوغ عند $e_n$ المعظِّمة: $\vertiii T = \max s_n$. وببتر
تفكيك القيم المفردة عند الرتبة $N$ يبقى مؤثر معياره $\sup_{n>N}s_n \to 0$: أي
التقريب المنتهي الرتبة. وليست لمؤثر فولتيرا أي قيمة ذاتية
(\cref{exo:b3:spectral:7})، لكن للمؤثر $V^*V$ قيمًا ذاتية: إذ $V^*Vf(x) = \int_x^1\int_0^tf(s)\,\dd s\,\dd t$ مؤثر
نواة متناظر موجب (بالنواة $1 - \max(x,y)$، وهي نواة غرين من
نمط الوتر)، وأزواجه الذاتية --- المحسوبة عبر مسألة الشروط
الحدّية $-u'' = \lambda^{-1}u$ مع $u'(0)
= u(1) = 0$، أي عائلة \cref{exo:b3:spectral:9} --- تعطي القيم
المفردة $s_n = \bigl((n + \frac12)\pi\bigr)^{-1}$. ويعيش تفكيك القيم المفردة على عائلتين
متعامدتين متجانستين \emph{اثنتين} بالضبط لأن $V$ يدير هندسته
الذاتية بعيدًا: فلا متجهات ذاتية، ومع ذلك بنيةٌ قطرية تامة بين
أساسين مختلفين.
\end{solution}

\begin{solution}{pb:b3:spectral:1}
\textbf{1.} الاتصال: لأن $\min$ و $\max$ متصلتان؛ والتناظر:
لأن تبديل $x, y$ لا يغيّر $\min(x,y)$ ولا $1 -
\max(x,y)$. والحدود:
$0 \leq g$، وتعطي حدود من نمط $g(x,y) \leq
\max\cdot(1-\max)$ أن $g \leq \frac14$ (فمن
أجل $u
= \max$: $\min \leq u$ ومنه $g \leq u(1-u) \leq \frac14$). و $g
\in L^2(\square)$: أي
هيلبرت--شميدت، ومنه $G$ متراص (\cref{exo:b3:spectral:4})؛ والنواة حقيقية
متناظرة: ومنه $G$ \omterm{def:b3:spectral:selfadjoint}{ذاتي المرافقة}.

\textbf{2.} بالشطر عند $y = x$:
\[
u(x) = (1 - x)\int_0^x y\,f(y)\,\dd y +
x\int_x^1(1 - y)\,f(y)\,\dd y .
\]
ومن أجل $f$ متصلة، نشتقّ (بقاعدة الجداء والمبرهنة الأساسية):
\[
u'(x) = -\int_0^xyf + \int_x^1(1-y)f
\qquad\text{(إذ تتلاشى الحدود الحدّية)},
\]
و $u''(x) = -xf(x) - (1 - x)f(x) = -f(x)$؛ ومن الواضح أن $u(0) =
u(1) = 0$. وبالعكس إذا كانت $u \in \mathcal C^2$
منعدمة عند الطرفين، فإن $w = u - G(-u'')$ يحقق $w'' = 0$ مع $w(0) = w(1) =
0$: ومنه
$w$ تآلفي وينعدم مرتين، فيكون $w = 0$.

\textbf{3.} ليكن $Gf = 0$ مع $f \in L^2$. من أجل $\psi \in \mathcal
C_c^\infty(\intoo01)$:
$\psi = G(-\psi'')$ حسب السؤال 2، ومنه
\[
\langle f, \psi\rangle = \langle f, G(-\psi'')\rangle
= \langle Gf, -\psi''\rangle = 0
\]
(لأن $G$ \omterm{def:b3:spectral:selfadjoint}{ذاتي المرافقة}): وحسب المبرهنة المساعدة الأساسية
(\cref{cor:b3:lp:fundlemma})، يكون $f = 0$ في كل مكان تقريبًا. وأما الإيجابية:
فمن أجل $f$ متصلة، مع $u = Gf$،
\[
\langle f, Gf\rangle = \int_0^1 fu = \int_0^1(-u'')u
= \bigl[-u'u\bigr]_0^1 + \int_0^1(u')^2 = \int_0^1(u')^2
\geq 0 ;
\]
ومن أجل $f \in L^2$، نقرّب في $L^2$ بدوال متصلة $f_n$: فينتقل
الطرفان إلى النهاية (لأن $G$ محدود).

\textbf{4.} إذا كان $Gu = \lambda u$ مع $\lambda \neq 0$: فإن
$Gu$ متصلة ($\abs{Gu(x) - Gu(x')} \leq \norm{g(x,\cdot) -
g(x',\cdot)}_2\norm u_2$، والنواة متصلة بانتظام)، ومنه تقبل $u$
ممثّلًا \omterm{def:b3:topology:continuity}{متصلًا}؛ عندئذٍ تبيّن صيغ السؤال 2 أن $Gu \in \mathcal C^2$، ومنه
$u =
\frac1\lambda Gu \in \mathcal C^2$ مع $-\lambda u'' =
-(Gu)'' = u$ و $u(0) = u(1) = 0$.

\textbf{5.} $-\lambda u'' = u$ مع $u(0) = 0$: $u = A\sin(x/\sqrt
\lambda)$ (و $\lambda$ موجب:
حسب السؤال 3، $\lambda =
\langle u, Gu\rangle/\norm u^2 > 0$ على المتجهات الذاتية). ويفرض $u(1) =
0$ أن
$\frac1{\sqrt\lambda} = n\pi$: أي $\lambda_n =
\frac1{n^2\pi^2}$، بدوال ذاتية $\sin(n\pi x)$، معيَّرة
$e_n = \sqrt2\sin(n\pi x)$ ($\int_0^12\sin^2(n\pi x)\dd x =
1$). وفحص التعامد: $2\sin(m\pi x)\sin(n\pi x) =
\cos((m-n)\pi x) - \cos((m+n)\pi x)$ يكامَل إلى $0$ من أجل
$m
\neq n$.

\textbf{6.} $\ker G = \{0\}$ (السؤال 3)، ومنه تعطي
\cref{thm:b3:spectral:spectral}(1) أن $H =
\overline{\operatorname{Vect}}(e_n)$: أي إن الجيوب أساسٌ هيلبرتي للفضاء
$L^2(\intcc01)$. ومن أجل $f(x) = x(1 - x)$:
\[
c_n = \sqrt2\int_0^1x(1-x)\sin(n\pi x)\,\dd x
= \sqrt2\;\frac{2\bigl(1 - (-1)^n\bigr)}{n^3\pi^3}
= \begin{cases}\dfrac{4\sqrt2}{n^3\pi^3} & n \text{ فردي},\\
0 & n \text{ زوجي},\end{cases}
\]
(بمكاملتين بالتجزئة). وببارسيفال: $\int_0^1x^2(1-x)^2\dd
x = \frac1{30} = \sum_{n \text{ فردي}}\frac{32}{n^6\pi^6}$، أي $\sum_{n\text{ فردي}}n^{-6} = \frac{\pi^6}{960}$.

\textbf{7.} $\langle e_n, Ge_n\rangle = \lambda_n\norm{e_n}^2
= \lambda_n$. ومن أجل $x$ مثبَّت، تكون معاملات $g(x,
\cdot)$
هي $\langle e_n, g(x,\cdot)\rangle = (Ge_n)(x) =
\lambda_ne_n(x)$، ومنه $g(x,\cdot) =
\sum_n\lambda_ne_n(x)\,e_n$ في $L^2$. والمتسلسلة الصريحة $\sum_n\lambda_ne_n(x)e_n(y) = \sum_n\frac{2\sin(n\pi
x)\sin(n\pi y)}{n^2\pi^2}$
تتقارب تقاربًا ناظميًّا على المربّع ($\abs{\text{الحدّ}} \leq \frac2{n^2\pi^2}$): فيكون مجموعها
\omterm{def:b3:topology:continuity}{متصلًا}، وله من أجل كل $x$ معاملات $L^2(\dd y)$ نفسها التي
للدالة $g(x, \cdot)$: ومنه تتوافق الدالتان المتصلتان من أجل كل
$(x, y)$. وبوضع $y = x$ والمكاملة (إذ يسمح التقارب الناظمي
بالمكاملة حدًّا حدًّا):
\[
\int_0^1g(x,x)\,\dd x =
\sum_n\lambda_n\int_0^1e_n(x)^2\dd x = \sum_n\lambda_n .
\]

\textbf{8.} $\int_0^1g(x,x)\dd x = \int_0^1x(1 - x)\dd x =
\frac16$، ومنه $\sum_{n\geq1}\frac1{n^2\pi^2} = \frac16$:
\[
\zeta(2) = \sum_{n\geq1}\frac1{n^2} = \frac{\pi^2}6 .
\]

\textbf{9.} من أجل $f = \mathbf 1$: $c_n =
\sqrt2\int_0^1\sin(n\pi x)\dd x = \sqrt2\,\frac{1 -
(-1)^n}{n\pi}$: أي $c_n = \frac{2\sqrt2}{n\pi}$ من أجل
$n$ فردي، و $0$ من أجل الزوجي. وببارسيفال: $1 = \sum_{n\text{ فردي}}\frac{8}{n^2\pi^2}$، ومنه
$\sum_{n\text{ فردي}}n^{-2} =
\frac{\pi^2}8$، و $\zeta(2) = \frac{\pi^2}8\cdot\frac{1}{1
- \frac14} = \frac{\pi^2}6$ (لأن الحدود الزوجية هي $\frac14\zeta(2)$). وأما
المقارنة: \omterm{thm:b3:hilbert:parseval}{فبارسيفال} من أجل $f$ واحدة تجمع $\abs{\langle e_n, f\rangle}^2$؛ بينما تكامل
صيغة الأثر قطرَ \emph{النواة}، وهو يعادل جمع \omterm{thm:b3:hilbert:parseval}{بارسيفال} على عائلة
متعامدة متجانسة كاملة دفعةً واحدة --- $\sum_n\langle e_n, Ge_n\rangle$ --- ومنه فهي
عمياء عن أي اختيار خاص لدالة الاختبار.

\textbf{10.} $\abs{c_n\cos(n\pi t)} \leq \abs{c_n}$ مع $\sum\abs{c_n}^2 < \infty$: فمن أجل كل $t$ تتقارب
المتسلسلة في $L^2$ (بالنشر المتعامد المتجانس،
\cref{thm:b3:hilbert:parseval}(3))؛ وحدّ الذيل $\norm{u(t) - u_N(t)}_2^2 \leq \sum_{n>N}\abs{c_n}^2$ منتظم في $t$، وكل مجموع
جزئي متصل في $t$ (لأنه عدد منتهٍ من الجيبتماميات): ومنه $t \mapsto u(t,\cdot)$
متصلة بقيم في $L^2$. ومن أجل $f = \sum_{n\leq N}c_ne_n$: يحقق كل نمط $\cos(n\pi t)\sin(n\pi x)$
المعادلةَ $\partial_t^2 =
-n^2\pi^2 = \partial_x^2$ مطبَّقةً عليه، وينعدم عند $x = 0,
1$، وقيمته
$\sin(n\pi x)$ ومشتقته الزمنية $0$ عند $t =
0$: فيحل المجموع
المنتهي كل شيء. وموسيقيًّا: تكون حركة الوتر تراكبًا من \emph{أمواج
قائمة} $e_n$ تردداتها $n\pi$ هي التردد الأساسي وتوافقياته؛
وتقول المبرهنة الطيفية إن كل شكل ابتدائي يتفكّك بكيفية وحيدة
إلى هذه النغمات الصافية، وتكون المعاملات $c_n$ هي طابع الصوت.
فسماع وتر هو حسابُ نشر متعامد متجانس.

\textbf{11.} مع $a_n = \abs{\langle u_n, x\rangle}^2$: $m_{p+1} = \sum_n\mu_n^{p+1}a_n = \sum_n\bigl(\mu_n^{p/2}
\sqrt{a_n}\bigr)\bigl(\mu_n^{p/2+1}\sqrt{a_n}\bigr) \leq
\sqrt{m_p\,m_{p+2}}$ (بكوشي--شوارتز في $\ell^2$).
ومنه تكون النسب $m_{p+1}/m_p$ متزايدة بالمعنى الواسع في $p$؛
وبما أن $R(x) = \frac{m_1}{m_0}$ و $\frac{\langle Ax, Ax\rangle}
{\langle x, Ax\rangle} = \frac{m_2}{m_1}$ و $R(Ax) =
\frac{m_3}{m_2}$، تنتج السلسلة، وكل حدّ منها
$\leq \mu_1$ لأن $m_{p+1} \leq \mu_1m_p$ حدًّا حدًّا. وأما التقارب: فإذا كان
$a_1 > 0$ (بكتابة الوزن الكلي للقيمة الذاتية العليا $a_1$)،
فإن
\[
\mu_1 \geq R(A^kx) = \frac{m_{2k+1}}{m_{2k}} =
\mu_1\,\frac{a_1 + \sum_{\mu_n<\mu_1}(\mu_n/\mu_1)^{2k+1}
a_n}{a_1 + \sum_{\mu_n<\mu_1}(\mu_n/\mu_1)^{2k}a_n}
\longrightarrow \mu_1,
\]
بالتقارب المهيمن للمجاميع (إذ النسب $< 1$): فتتقارب طريقة القوى
من أجل كل متجهة انطلاق غير متعامدة مع الفضاء الذاتي الأعلى.

\textbf{12.} نقطةً نقطة، $\abs{\langle x, (A - B)x\rangle}
\leq \vertiii{A - B}\,\norm x^2$، ومنه $R_A(x) \leq R_B(x) +
\vertiii{A - B}$ من أجل كل $x$.
وبإدخال ذلك في صيغة الأعظم--الأصغر في \cref{exo:b3:spectral:8}: $\mu_n(A)
\leq \mu_n(B) + \vertiii{A - B}$،
وبالتناظر في $A,
B$: $\abs{\mu_n(A) - \mu_n(B)} \leq \vertiii{A - B}$ من أجل كل $n$ دفعةً واحدة.

\textbf{13.} يُكامل $-w'' = x - x^2$ ليعطي $w = -\frac{x^3}6
+ \frac{x^4}{12} + cx$ (مع
$w(0) = 0$)، ويعطي $w(1) = 0$ أن $c = \frac1{12}$:
\[
w = \frac{x^4 - 2x^3 + x}{12} = \frac{x(1-x)(1 + x -
x^2)}{12} = Gu .
\]
ثم $\norm u_2^2 = \int_0^1x^2(1-x)^2 = \frac1{30}$، ومع $\int_0^1x^3(1-x)^3 = B(4,4) = \frac1{140}$:
\[
\langle u, Gu\rangle = \frac1{12}\Bigl(\frac1{30} +
\frac1{140}\Bigr) = \frac{17}{5040},
\qquad
R(u) = \frac{17/5040}{1/30} = \frac{17}{168} .
\]
ومنه $\frac1{\pi^2} = \lambda_1 \geq \frac{17}{168}$، أي $\pi^2 \leq \frac{168}{17} = 9.8824$: أي $\pi \leq 3.14364 <
3.1437$ (والقيمة الحقيقية
$\pi^2 = 9.8696$). كثير حدود، وتكامل، ورقم.

\textbf{14.} نكتب $u = x - x^2$، ومنه $Gu =
\frac{u(1 + u)}{12}$، وباستعمال
$\int u^2 = \frac1{30}$ و $\int u^3 = \frac1{140}$ و $\int u^4 = B(5,5) =
\frac{4!\,4!}{9!} = \frac1{630}$:
\[
\norm{Gu}_2^2 = \frac1{144}\int u^2(1+u)^2 =
\frac1{144}\Bigl(\frac1{30} + \frac2{140} +
\frac1{630}\Bigr) = \frac1{144}\cdot\frac{62}{1260} =
\frac{31}{90720} .
\]
ومنه $\frac{\langle Gu, Gu\rangle}{\langle u, Gu\rangle} =
\frac{31/90720}{17/5040} = \frac{31}{306}$، وحسب السؤال 11 يبقى هذا $\leq \mu_1 = \frac1{\pi^2}$: أي $\pi^2 \leq
\frac{306}{31} = 9.87097$، أي
$\pi \leq 3.14181 <
3.1419$ --- أربعة أرقام (والتكرار التالي كان سيعطي نحو ثمانية،
إذ ينكمش الخطأ بالعامل $(\lambda_2/\lambda_1)^2 = \frac1{16}$ في كل خطوة).

\textbf{15.} العائلة $(e_m \otimes e_n)(x,y) =
e_m(x)e_n(y)$ أساسٌ هيلبرتي للفضاء $L^2(\intcc01^2)$
(بالتعامد والتجانس بتونيلي؛ وبالشمول كما في \cref{exo:b3:spectral:5}).
وحسب السؤال 7، من أجل $x$ مثبَّت: $g(x, \cdot) = \sum_n\lambda_ne_n(x)e_n$، ومنه يكون معامل $g$
على $e_m\otimes e_n$ هو
\[
\langle e_m\otimes e_n,\ g\rangle
= \int_0^1 e_m(x)\,\lambda_n e_n(x)\,\dd x
= \lambda_n\,\delta_{mn} .
\]
وببارسيفال في المربّع:
\[
\iint g^2 = \sum_{m,n}\abs{\langle e_m\otimes e_n,
g\rangle}^2 = \sum_n\lambda_n^2 .
\]

\textbf{16.} بتناظر $g$،
\[
\iint g^2 = 2\iint_{y<x}y^2(1-x)^2
= 2\int_0^1(1-x)^2\,\frac{x^3}3\,\dd x
= \frac23\,B(4, 3) = \frac23\cdot\frac{3!\,2!}{6!}
= \frac23\cdot\frac1{60} = \frac1{90} .
\]

\textbf{17.} وبالضمّ: $\sum_n\frac1{n^4\pi^4} =
\frac1{90}$، أي $\zeta(4) = \frac{\pi^4}{90}$. وعمومًا، يساوي
$\operatorname{tr}(G^k) = \sum\lambda_n^k =
\frac{\zeta(2k)}{\pi^{2k}}$ تكاملًا متتاليًا لجداءات النواة الكسرية الحدودية $g$
على المكعب ذي البُعد $k$: أي عددًا ناطقًا. ومن ثَمّ $\zeta(2k) \in
\pi^{2k}\Q$ من
أجل كل $k$. ولا تبلغ الآلة إلا الوسائط الزوجية لأن القيم الذاتية
تدخل عبر \emph{قواها} --- $\sum\lambda_n^k$ --- و $\lambda_n =
\frac1{n^2\pi^2}$: فلا تنتج أي
تركيبة من الآثار المقدارَ $\sum
n^{-3}$؛ وتقع الطبيعة الحسابية للمقدار
$\zeta(3)$ (وهو غير ناطق حسب أبيري، وتساميه مفتوح) وراء
المحاسبة الطيفية.

\textbf{18.} الحدّ الأعلى لمجموع حدود موجبة أصغر من أو يساوي
المجموع: $\lambda_1^2 \leq \sum\lambda_n^2 =
\frac1{90}$، ومنه $\frac1{\pi^2} \leq \frac1{\sqrt{90}}$ و $\pi \geq 90^{1/4} = 3.0801\ldots$ ومع السؤال 14:
$3.080 < \pi < 3.1419$، بحساب الوتر وحده. وتحدّد الآثار العليا الحدَّ الأدنى
هندسيًّا: $\lambda_1 \leq (\operatorname{tr}G^{2k})^{1/2k} =
\lambda_1\bigl(1 + \sum_{n\geq2}(\lambda_n/\lambda_1)^{2k}
\bigr)^{1/2k}$، ويموت العامل الطفيلي بسرعة $\bigl(\tfrac14\bigr)^{2k}\cdot\frac1{2k}$ --- وهي
آلية الفجوة الطيفية نفسها التي يقوم عليها تقارب طريقة القوى
(السؤال 11)، منظورًا إليها من جهة الأثر.

\textbf{19.} $\abs{n^2\pi^2 - \nu} \geq \delta > 0$ من أجل كل $n$ (إذ تتفادى المتتالية $n^2\pi^2 \to \infty$
القيمةَ $\nu$ بهامش)، و $\abs{n^2\pi^2 - \nu} \geq \frac{n^2\pi^2}2$ من أجل $n$ كبير. وأما التقارب
في $L^2$: فالمعاملات $\frac{c_n}{n^2\pi^2 - \nu}$ قابلة للجمع بالمربّع (مهيمَنًا
عليها بالمقدار $\frac{\abs{c_n}}\delta$). وأما التقارب المنتظم: فتُحدّ نظائم
سوپريموم الذيول بالمقدار $\sqrt2\sum_{n>N}
\frac{\abs{c_n}}{\abs{n^2\pi^2 - \nu}} \leq
\frac{2\sqrt2}{\pi^2}\bigl(\sum\abs{c_n}^2\bigr)^{1/2}
\bigl(\sum_{n>N}n^{-4}\bigr)^{1/2} \to 0$ (بكوشي--شوارتز). وأما التحقق:
فللمقدار $Gf + \nu Gu$ المعاملُ على $e_n$
\[
\lambda_nc_n + \frac{\nu\lambda_nc_n}{n^2\pi^2 - \nu}
= \frac{c_n}{n^2\pi^2}\Bigl(1 + \frac{\nu}{n^2\pi^2 -
\nu}\Bigr) = \frac{c_n}{n^2\pi^2 - \nu} :
\]
وهي بالضبط معاملات $u$، ومنه $u = Gf + \nu Gu$؛ وأما الوحدانية فلأن فرق
حلّين $v$ يحقق $v = \nu
Gv$، أي $\langle e_n, v\rangle(n^2\pi^2 - \nu) = 0$ من أجل كل $n$: ومنه $v = 0$.

\textbf{20.} كما في السؤال 19، تكافئ المعادلة $u = Gf + \nu
Gu$ عائلةَ
معادلات المعاملات $(n^2\pi^2 - \nu)\,\langle e_n, u\rangle = c_n$، $n \geq
1$. ومن أجل $n = m$ يكون الطرف
الأيسر $0$: فتفرض قابلية الحل أن $c_m = 0$، ويكون عندئذٍ $\langle e_m,
u\rangle$
حرًّا بينما تتعيّن جميع المعاملات الأخرى: فتشكّل الحلول المستقيم
$u_0 + \R e_m$. وأما الرنين: فإقحامٌ ذو مركّبة على النمط الذاتي
يضخّ فيه الطاقة بلا حدّ --- أي الأرجوحة المدفوعة عند ترددها
الخاص.

\textbf{21.} من صيغة السؤال 19، $\norm{R_\nu f}_2^2
= \sum\frac{\abs{c_n}^2}{(n^2\pi^2 - \nu)^2} \leq
\frac{\norm f^2}{(\pi^2 - \nu)^2}$ (إذ من أجل $\nu < \pi^2$
تكون القيمة الذاتية الأقرب هي $\pi^2$)، مع مقاربة التساوي عند
$f = e_1$: أي \omterm{def:b3:banach:operator}{معيار المؤثر} $\frac1{\pi^2 - \nu}$. وأما التراص: فإن $R_\nu$
نهايةٌ بالمعيار لبتوره المنتهية الرتبة (إذ معاملات الذيل $\frac1{n^2\pi^2 - \nu} \to 0$)؛
وتُقرأ ذاتية المرافقة والإيجابية من الصورة القطرية (إذ جميع
المعاملات $\frac1{n^2\pi^2 - \nu} > 0$). وللمؤثر $R_\nu$ القيمُ الذاتية $\frac1{n^2\pi^2 - \nu}$:
فيبدأ تحليل الأجزاء من الأول إلى السادس من جديد حرفيًّا.

\textbf{22.} المعجم: \emph{القيمة الذاتية} $\lambda_n =
\frac1{n^2\pi^2}$ $\leftrightarrow$ مربّع
التردد $n^2\pi^2$ للتوافقي ذي الرتبة $n$؛ و\emph{الأثر} $\sum
\lambda_n = \frac16$
$\leftrightarrow$ $\zeta(2) =
\frac{\pi^2}6$؛ و\emph{معيار هيلبرت--شميدت} $\iint g^2 =
\frac1{90}$ $\leftrightarrow$
$\zeta(4) = \frac{\pi^4}{90}$؛ و\emph{\omterm{thm:b3:spectral:fredholm}{بديل فريدهولم}} $\leftrightarrow$ رنين الوتر المُقحَم؛
و\emph{الأصغري--الأعظمي} $\leftrightarrow$ التقديرات التغيّرية، وصولًا
إلى $\pi < 3.1437$ من كثير حدود واحد. ووراء كل اقتران، الكائن
نفسه: \omterm{def:b3:spectral:compact}{مؤثر متراص} \omterm{def:b3:spectral:selfadjoint}{ذاتي المرافقة} واحد، مقطرن مرة واحدة، مستثمَر
خمس مرات.

\textbf{23.} بالشطر على التماثل وبالتعويض $n = 2m$ في الجزء
الزوجي،
\[
\zeta(6) = \sum_{n\text{ فردي}}\frac1{n^6} +
\sum_{m\geq1}\frac1{(2m)^6}
= \frac{\pi^6}{960} + \frac{\zeta(6)}{64},
\]
ومنه $\frac{63}{64}\zeta(6) = \frac{\pi^6}{960}$ و $\zeta(6)
= \frac{64\,\pi^6}{63\cdot960} = \frac{\pi^6}{945}$. وكان طريق الأثر سيحسب $\operatorname{tr}G^3 =
\sum\lambda_n^3 = \zeta(6)/\pi^6$ بوصفه
$\iint g\,g_2$ بالنواة المكرَّرة $g_2(x,y) = \int_0^1g(x,z)g(z,y)\dd z$ --- أي ثلاث مكاملات
لكثيرات حدود بالقطع؛ بينما لم تحتج \omterm{thm:b3:hilbert:parseval}{بارسيفال} على $x(1-x)$ إلا
إلى واحدة.

\textbf{24.} (a) نقطرن: $v = \sum_na_nu_n$ (مع مركّبة نواة ممكنة، ولا يكسب
عليها $\langle v, Av\rangle$ شيئًا بينما ينمو $\norm v^2$، ومنه لا تكون
لمعظِّم أي مركّبة منها). عندئذٍ $\langle v, Av\rangle = \sum\mu_na_n^2 \leq
\mu_1\sum a_n^2$، مع التساوي إذا وفقط
إذا كان $a_n = 0$ كلما كان $\mu_n < \mu_1$: أي إن المعظِّم يقع
في الفضاء الذاتي الموافق للقيمة $\mu_1$. (b) ومن أجل أي $u$،
\[
\langle\abs u, A\abs u\rangle - \langle u, Au\rangle
= \iint k(x,y)\,\bigl(\abs{u(x)}\abs{u(y)} -
u(x)u(y)\bigr)\dd x\,\dd y \;\geq\; 0,
\]
لأن المقدار المكامَل غير سالب نقطةً نقطة. وإذا كان لكلٍّ من
$P = \{u >
0\}$ و $N = \{u < 0\}$ قياسٌ موجب، فإن المقدار المكامَل يساوي
على $P\times N$ المقدارَ $2k\abs{u(x)}\abs{u(y)} >
0$ على مجموعة ذات \omterm{def:b3:measure:measure}{قياس} موجب: أي
متراجحة أكيدة. ودالةٌ ذاتية $u$ للقيمة $\mu_1$ تعظّم نسبة رايلي،
وللمقدار $\abs u$ المعيار نفسه، ومنه كانت الأكيدية ستُبرز $R(\abs u) > \mu_1$
--- وهو مستحيل؛ فتكون للدالة $u$ إشارة ثابتة في كل مكان تقريبًا،
ولنقل $u \geq 0$. عندئذٍ $u(x) = \mu_1^{-1}(Au)(x) =
\mu_1^{-1}\int k(x,y)u(y)\dd y > 0$ من أجل كل $x \in
\intoo01$ (لأن
$k(x,\cdot) > 0$ و $u \neq 0$). (c) ولو كان بُعد الفضاء الذاتي
$\geq 2$، لاحتوى دالتين ذاتيتين \emph{متعامدتين} $u, v$، لكلٍّ
منهما إشارة ثابتة ولا تنعدم داخليًّا حسب (b)؛ لكن عندئذٍ $\abs{\langle u, v\rangle} = \int\abs u\,\abs v > 0$
--- وهو تناقض. وعلى الوتر: $k = g > 0$ على المربّع المفتوح،
و $\mu_1 = \lambda_1 = \frac1{\pi^2}$ بسيطة فعلًا، و $e_1 = \sqrt2\sin(\pi x)$ موجبة على $\intoo01$؛ وكل
$e_n = \sqrt2\sin(n\pi x)$ حيث $n \geq 2$، بكونها متعامدة مع $e_1$ الموجبة، يجب أن
يكون تكاملها مقابلها معدومًا، ومنه تغيّر إشارتها --- كما تؤكد
جذورها الداخلية $n - 1$ عند $\frac kn$.

\textbf{25.} بما أن $n^2\pi^2 \to \infty$، تُبلَغ الأصغرية $d =
\min_n\abs{n^2\pi^2 - \nu}$ عند نمط
$m$ ما، ويكون $d > 0$ لأن $\nu$ يتفادى الطيف. وتعطي صيغة السؤال
19 القطرية أن
\[
\norm{R_\nu f}_2^2 =
\sum_n\frac{\abs{c_n}^2}{(n^2\pi^2 - \nu)^2}
\leq \frac1{d^2}\,\norm f_2^2,
\]
مع التساوي من أجل $f = e_m$: أي $\vertiii{R_\nu} = \frac1d$، وهو مقلوب المسافة من
$\nu$ إلى الطيف --- أي مبدأ المُحلِّل العام، هنا بالإحداثيات
الصريحة. وتُقرأ ذاتية المرافقة من المعاملات القطرية الحقيقية؛
وينتج التراص كما في السؤال 21 (إذ تؤول المعاملات إلى $0$، ومنه
تتقارب البتور المنتهية الرتبة بالمعيار). وأما ثمن الرنين: فمن
أجل $f = e_1$ و $\nu =
(1-\varepsilon)\pi^2$، تعطي الصيغة $u =
\frac{c_1}{\pi^2 - \nu}e_1 = \frac1{\varepsilon\pi^2}e_1$، مقابل الاستجابة
الساكنة $Ge_1 = \frac1{\pi^2}e_1$: أي تضخيم $\frac1\varepsilon$. وعند $\varepsilon =
10^{-2}$ تكون الاستجابة
$100$ ضعف الساكنة --- وهي تتباعد حين $\varepsilon \to 0$، وذلك هو بديل
السؤال 20 منظورًا إليه من الجهة المحدودة: فكلما اقترب تردد
الإقحام من تردد طبيعي، قلّ حدّ المقلوب.
\end{solution}
